\section{PX4 HITL simulation and validation}
\label{sec:test-4-hitl}

% Setup:    build drone, HITL configuration
% Test:     - AirSim + follow on Windows + Pixhawk
%           - QGroundControl (without program running)
%           - PX4 to computer connection
%           - RC on AirSim
% Results:  images of wiring computer-Pixhawk

The purpose of this section is to validate the transition from using a simulated version of the flight stack running on Linux (PX4's software-in-the-loop simulation) to using a physical Pixhawk board with simulated input and output to test the flight controller interaction with the developed program.
To do so, the aim is to be able to run the follow solution to send control commands through the Pixhawk board and observe the movement of the vehicle in the AirSim simulator in the same manner as in the previous section.


To activate this new hardware-in-the-loop mode, QGroundControl contains a specific quadcopter HITL airframe configuration that sets up the board with all the required parameters.
QGroundControl automatically detects the Pixhawk 4 board when it is connected to the computer through its Micro-USB port.
It is as well required to make changes to the AirSim configuration for the simulator to work with HITL mode, namely, activating the option to accept connections through serial.
To test the complete system configuration for HITL described in section \ref{sec:devenv} and outlined in figure \ref{fig:hitl-connections} the Pixhawk board needs an additional channel of communication to the computer dedicated to the Mavlink exchange with the dronecontrol application.
The board will therefore have both a direct cable connection to the computer and a telemetry radio on its \verb|TELEM1| port with a wireless link to its counterparty radio connected to the computer.
Since the AirSim simulator requires a higher update rate than the dronecontrol application it will employ the faster, cabled link, through which it can automatically connect to the board when it is started.
The dronecontrol program will connect through the telemetry radio by specifying the serial port and its baudrate.
Figure \ref{fig:hitl-setup-picture} shows all the connections mentioned.

\begin{figure}
  \centering
  \includegraphics[width=.6\textwidth, keepaspectratio]{img/placeholder.png}
  \caption{Pixhawk 4 board connected to the computer running the AirSim simulator and the dronecontrol application}\label{fig:hitl-setup-picture}
\end{figure}

After all the necessary connections are set up, the program can be started with the following command:
\mintinline{bash}{python -m dronecontrol follow --sim --serial COM[X]:57600}, where the exact COM port will vary depending on the particular USB port to which the telemetry radio is connected.


With the results matching exactly those obtained in section \ref{sec:test-3-airsim}, and since the flight stack is now running on a physical controller, it is possible to add an RC antenna to the \verb|PPM RC| \todo{check port} port of the board to be able to fly the vehicle with a remote control unit after it has been bound to the receiver \footnote{\url{https://docs.px4.io/main/en/config/radio.html}}.
By configuring one of the switches in the RC unit to change to PX4's offboard flight mode, additional checks to the safety measures of interrupting autonomous flight on flight mode changes or loss of signal from the RC controller can be carried out.
As can be seen in figure \ref{fig:rc-tests}, ....

\begin{figure}
  \centering
  \includegraphics[width=.6\textwidth, keepaspectratio]{img/placeholder.png}
  \caption{Some figures that show these tests}\label{fig:rc-tests}
\end{figure}


\subsection{PX4 HITL validation with Raspberry Pi}
\label{sec:test-5-rpi}

% Setup:    Raspberry Pi installation, RPi-Pixhawk connection
% Test:     - AirSim + follow on RPi + Pixhawk
%           - Serial connection
% Results:  wiring, performance metrics

The next step in the transition from a fully simulated environment to real flight is to connect the future onboard computer, the Raspberry Pi 4, to the Pixhawk flight controller and proceed with more realistic tests on the exact hardware that will be controlling the drone.
The main characteristics of the Raspberry Pi that have to be ensured to be able to progress further towards autonomous flight are:
\begin{enumerate}
    \item Capacity to function when power is provided from a battery
    \item Stability of serial connection to the Pixhawk board
    \item Ability to run the dronecontrol application and all its dependencies
    \item Connection to an external camera
    \item Performance of the computer vision algorithms with reduced processing power
\end{enumerate}

% installation
% xrdp
% battery tests

The complete installation process of all the required libraries and dependencies for the Raspberry Pi is explained in appendix \ref{app:install-dronecontrol-rpi}.
The most convenient method of controlling the small computer is through a remote desktop connection, where the screen contents and mouse and keyboard input are transmitted through a local network.
This way, it is possible to access the Pi's desktop from the ground station computer even during flight.
One option to achieve this is XRDP\footnote{\url{http://xrdp.org/}}, an open-source implementation of a Microsoft Remote Desktop Protocol server compatible with the Raspberry OS.

The first hardware connection to test is the power supply to the Raspberry.
The most common way of providing current to the Pi when mobility is not required is through its USB-C port directly to a wall socket.
This provides the recommended 3 amps that it needs to function at maximum capacity.
The power module included with the Holybros' X500 development kit has two ports that provide 5 volts with an output current of 3 amps.
One of them will therefore be used to power the flight controller and the other to power the Raspberry Pi, as described in section \ref{subsec:onboard}.
Figure \ref{fig:power-supply} shows the power supply through the battery to the two computers.

\todo[inline]{Measure voltage and amperage through power ports separately and together}
\todo[inline]{Close up image and description of custom connector}

\begin{figure}
  \centering
  \includegraphics[width=.6\textwidth, keepaspectratio]{img/placeholder.png}
  \caption{Picture of Rpi + power module + battery + pixhawk}\label{fig:power-supply}
\end{figure}


% serial activation (enable MAVLINK 2 with baudrate)

A similar connector is used to provide a channel for the Mavlink communication between the flight controller and the onboard computer.
In this case, one end of the connector is attached to the \verb|TELEM2| port in the board and the TX/RX pins are attached to the UART pins on the Raspberry's GPIO header.
For this serial connection to work, both the flight controller and the companion computer need some additional configuration.
The Pixhawk needs to be configured through QGroundControl to enable a secondary Mavlink channel as, by default, only the \verb|TELEM1| port used by the telemetry radio is configured.
The necessary parameters to modify and their values are collected in table \ref{tab:telem2-params}.
On the Raspberry side, the serial port is configured to be used as a terminal by default.
This can be disabled through the \verb|raspi-config| command-line utility (Advanced config -> Serial -> Disable\todo{Check option path on latest OS}).
After that, the \verb|/dev/serial0| address can be used to communicate to the device attached to the UART pins at the baudrate configured through QGroundControl.

\todo[inline]{Fix table}
\begin{table}[h!]
 \begin{center}
  \begin{tabular}{l|l}
    Parameter name & Value \\ \hline
    MAV\_2\_CONFIG & TELEM2 \\
    SER\_TEL2\_BAUD & 921600 \\
  \end{tabular}
  \caption{PX4 parameters than need to be configured to enable Mavlink communication through the secondary telemetry port}
  \label{tab:telem2-params}
 \end{center}
\end{table}


\begin{figure}
  \centering
  \includegraphics[width=.6\textwidth, keepaspectratio]{img/placeholder.png}
  \caption{a) Picture of Rpi raspi-config || b) close-up of pixhawk to pi cable connection}\label{fig:serial-connection}
\end{figure}

The test camera utility will be used to validate this configuration.
At this point, the only physical connection to the ground station running AirSim and/or QGroundControl is through the development-only micro-USB port in the Pixhawk.
The results from running: \\
\begin{listing}[h!]
    \begin{minted}[breaklines, fontsize=\footnotesize, baselinestretch=1]{bash}
dronecontrol tools test-camera --hardware /dev/serial0:921600 --sim <AirSim host IP> --pose-detection
    \end{minted}
\end{listing}\\
can be seen on figure \ref{fig:rpi-airsim-test}.
On the left side, the remote connection to the Raspberry's desktop shows the output of the dronecontrol program running the pose detection algorithm on the images received from the simulator.
On the right side, the AirSim simulator shows the movements of the vehicle as it reacts to the input from the flight controller and the companion computer.

\begin{figure}
  \centering
  \includegraphics[width=.8\textwidth, keepaspectratio]{img/placeholder.png}
  \caption{a) RPi desktop with pose output || b) AirSim on Windows}\label{fig:rpi-airsim-test}
\end{figure}


% performance test ---> graphs
The main question to answer now is whether the less powerful processor in the Raspberry Pi 4b, a quad-core ARM Cortex-A72 64-bit SoC running at 1.5GHz, can handle the detection and tracking algorithms with a performance good enough to react to real-time movement.
\todo[inline]{Performance analysis with/out peripherals, detection only, detection and control together}

